\section{第~\ref{sec:debugging}~章　调试这台机器}

电极听到了什么、活动的低维性、基于刚性阵列的接口的工作、大脑与解码器的共同学习以及刺激引起的视觉光点，都已在同行评审的实验中测量；数据流的估计是本章的算术；关于植入人体的 Neuralink 设备，就我们所能核实的，数据主要由公司自己发布。

{\footnotesize
\setlength{\tabcolsep}{3pt}\renewcommand{\arraystretch}{1.15}
\begin{longtable}{>{\raggedright}p{0.5cm}>{\raggedright}p{4.0cm}>{\raggedright}p{5.6cm}>{\raggedright}p{2.3cm}>{\raggedright\arraybackslash}p{1.7cm}}
\toprule
编号 & 论断 & 实验及测量内容 & 来源 & 状态 \\
\midrule\endfirsthead
\toprule
编号 & 论断 & 实验及测量内容 & 来源 & 状态 \\
\midrule\endhead
1 & 电极能听到约一百微米半径内神经元的脉冲，通常是一个到几个；可按波形区分它们 & 大鼠 CA1 区，同时进行细胞内和细胞外记录：细胞外脉冲的波形重现了细胞内脉冲的宽度和幅度。估计：四极电极本可区分 $60$--$100$ 个神经元，实际只分出约六个。 & \cite{Henze2000extra} & 已测量 \\
2 & 大脑中有 $8.6\cdot10^{10}$ 个神经元；探针听到的约为一亿分之一 & 对溶解组织悬液中的细胞核计数：人脑中有 $86\cdot10^9$ 个神经元。比例是本章的算术。 & \cite{Azevedo2009} & 已测量 \\
3 & 运动皮层几百个神经元的活动可以放进一二十个公共变量 & 猴子运动皮层记录的综述：群体活动由少数为许多神经元所共有的隐变量描述。 & \cite{Gallego2017} & 已测量 \\
4 & 相邻神经元的噪声是相关的，一百个神经元携带的信息几乎和一千个一样多（命题~\ref{prop:pooling}） & 猴子视觉皮层：相邻神经元噪声的相关系数约为 $0.1$，平均化的收益在一百个神经元时饱和。限制信息的相关性理论。 & \cite{Zohary1994, MorenoBote2014} & 已测量 \\
5 & 原始数据流 $2\cdot10^8$~bit/s，而脉冲中只有 $8\cdot10^4$~bit/s~\eqref{eq:probedata} & 本章依据脉冲流容量~\eqref{eq:spikeentropy}~所作的算术。数字化参数是真实的：Neuralink 芯片为 $19.3$~kHz、每通道 $10$ 位。 & 本章推导；\cite{Rieke1997, Musk2019} & 理论 \\
6 & Neuralink 的第一个系统：芯片放大并数字化，原始数据流经电缆送出，脉冲由外部程序提取；据公司介绍，用于人体的设备在芯片上提取脉冲，采用封闭外壳、无线电以及无线供电 & 公司的论文：完整的原始数据流经 USB-C 电缆送出，脉冲由芯片外的程序实时提取；板上压缩、无线供电和无线传输被列为临床设备的计划。关于植入人体的设备，只有公司的介绍。必须在芯片上提取脉冲是本章由~\eqref{eq:probedata}~得出的结论。 & \cite{Musk2019}；Neuralink 报告，无同行评审论文 & 已测量（公司论文）；人体部分为公司报告 \\
7 & 最多 $3072$ 个电极，分布在 $96$ 根细丝上，每根 $32$ 个；细丝由机器人避开血管插入 & 公司的论文：$3072$ 个电极分布在 $96$ 根细丝上，每根 $32$ 个；机器人每分钟插入 $6$ 根细丝（$192$ 个电极）；长期植入阵列的大鼠中，多达 $70\,\%$ 的电极能听到脉冲。 & \cite{Musk2019} & 已测量（公司论文） \\
8 & 人体中约一千个通道分布在 $64$ 根细丝上；部分细丝移位；光标约 $10$~bit/s & 仅有公司的介绍。在 PubMed 中没有检索到报告人体结果的同行评审论文；这与本章正文中的保留说法一致。 & Neuralink 报告；无同行评审论文 & 已测量（公司报告） \\
9 & 基于一百个电极阵列的接口可以控制光标 & 瘫痪患者，初级运动皮层中 $96$ 个电极：受伤三年后，移动手臂的意图仍会改变脉冲；解码器提供“神经光标”（电子邮件、电视），并能控制假手和机械臂。 & \cite{Hochberg2006} & 已测量 \\
10 & 用“意念之手”写字：每分钟约 $90$ 个字符，低于 $8$~bit/s & 手部瘫痪患者尝试手写，使用循环神经网络解码器：实时每分钟 $90$ 个字符，准确率 $94.1\,\%$，自动校正后超过 $99\,\%$。以比特计的估计是本章的算术。 & \cite{Willett2021} & 已测量 \\
11 & 说话的尝试可以每分钟约 $60$ 个词的速度转换成文字 & 失去清晰言语能力的患者，皮层中植入阵列：每分钟 $62$ 个词；词汇量 $50$ 个词时词错误率 $9.1\,\%$，词汇量 $125\,000$ 个词时为 $23.8\,\%$。 & \cite{Willett2023} & 已测量 \\
12 & 接口只取出容量的一小部分：单个神经元每秒几十比特，一百个神经元每秒几千比特（命题~\ref{prop:probe}） & 上界~\eqref{eq:spikeentropy}~是理论；与第~8、10 行的比较是本章的推论。瓶颈在于低维性和对编码的无知而不在导线，这一点未经直接实验检验。 & \cite{Rieke1997} & 理论 \\
13 & 大脑和解码器相互学习 & 猴子控制光标；解码器在闭环中适应，神经元同时重新学习：准确率提高，技能得以保持，受自己手臂运动的干扰也更小。把学习速度拉开的建议是工程经验法则，本章没有单独的实验。 & \cite{Orsborn2014coad} & 已测量 \\
14 & 通过电极的电流不分类型地激发周围的神经元和导线 & 小鼠和猫的皮层，双光子钙成像：即使是弱电流，也会通过直接激发直径几十微米体积内的轴突，激活分散在远至数毫米处的稀疏神经元。也就是说，激发没有选择性，但也不是连成一片的。 & \cite{Histed2009stim} & 已测量 \\
15 & 刺激视觉皮层会点亮光点；同时点亮最多 $16$ 个光点可以认出一些字母和物体边界 & 盲人，视觉皮层中植入 $96$ 个电极 $6$ 个月：单个电极的阈值为 $66.8\pm36.5$~\textmu A；同时刺激电极组（最多 $16$ 个，这是刺激器的上限）会降低阈值并产生可分辨的图案，能认出一些字母和物体边界。 & \cite{Fernandez2021} & 已测量 \\
16 & Neuralink 的第二个方向是向视觉皮层送入信号的设备 & 只有公司关于研发的介绍；没有找到关于其工作的同行评审数据。 & Neuralink 报告 & 假说 \\
17 & 广播型神经递质的浓度可以用化学传感器在组织中测量 & 大鼠脑片，碳纤维快速循环伏安法：测量了血清素的释放和再摄取；该物质会扩散到突触之外。 & \cite{Bunin1998} & 已测量 \\
\bottomrule
\end{longtable}}
